Multi-agent large language model (LLM) systems routinely rely on
role-playing to divide complex tasks. However, it remains unclear whether
these assigned roles genuinely enforce a division of labor or inadvertently
alter other critical system dynamics. We investigate what roles change
beyond task allocation. By ablating role instructions for three open-weight
models on two reasoning benchmarks (Omni-MATH and SciBench), we isolate the
effect of roles while holding the agents, topology, call budget, and
ground-truth evaluator fixed. Deleting role instructions largely dismantles
functional specialization, yet first-pass accuracy does not degrade.
Instead, roles matter through the decision rights they allocate. A
role-structured Reviewer acts as a strict gatekeeper that often withholds
submissions, so the team receives less external corrective feedback. In
exploratory comparisons where both teams have equal access to feedback,
role-free teams recover from initial failures more often and
achieve higher final accuracy in every setting. The gap vanishes in a post
hoc comparison at equal numbers of evaluator verdicts. In a single exploratory run,
forcing the role-structured team to submit in every cycle raises its
recovery to the role-free level for gemma-4, suggesting that the gap stems
mainly from how often each team submits. We conclude that roles govern information flow and decision
rights as well as work allocation. Consequently, designing and evaluating
effective multi-agent organizations requires looking beyond task division
to explicitly manage feedback loops and decision policies.
